\documentclass[11pt,french]{article}

\usepackage[utf8]{inputenc}
\usepackage[T1]{fontenc}
\usepackage[french]{babel}
\usepackage{lmodern}
\usepackage{amsmath,amsfonts,amssymb}
\usepackage[top=2.5cm, bottom=2.5cm, left=2.5cm, right=2.5cm]{geometry}
\usepackage{booktabs}
\usepackage{array}
\usepackage{diagbox}
\usepackage{enumitem}
\usepackage{fancyhdr}
\usepackage[colorlinks=true, linkcolor=blue!70!black, urlcolor=blue!70!black]{hyperref}

\renewcommand{\baselinestretch}{1.15}
\setlength{\parindent}{0pt}
\setlength{\parskip}{0.55em}

\pagestyle{fancy}
\fancyhf{}
\fancyhead[L]{\footnotesize Université Laval \\ Département de mathématiques et de statistique}
\fancyhead[R]{\footnotesize STT-2902 -- Modélisation statistique \\ Julien Miron}
\fancyfoot[C]{\thepage}
\renewcommand{\headrulewidth}{0.5pt}
\setlength{\headheight}{28pt}

% Largeur de colonne centrée pour les tableaux
\newcolumntype{P}[1]{>{\centering\arraybackslash}p{#1}}

\newcommand{\remarque}[1]{\medskip\noindent\textit{\textbf{Remarque.}} #1\medskip}

\begin{document}

\begin{center}
  {\LARGE\bfseries Test d'indépendance}
\end{center}

\medskip
\begin{center}
\small{Le matériel de ce chapitre est tiré des notes de cours rédigées par Anne-Sophie Charest et basées en partie sur un document d'Emmanuelle Reny-Nolin.}
\end{center}

\bigskip
\tableofcontents

\newpage

%%============================================================
\section{Introduction}
%%============================================================

Les chapitres précédents portaient sur l'inférence pour une variable continue (moyenne,
variance) ou une variable catégorique à deux modalités (proportion). Ce chapitre traite de
l'association entre \emph{deux} variables catégoriques.

La question centrale est la suivante~: peut-on conclure, à partir des données observées,
que les deux variables sont associées dans la population, ou la relation observée est-elle
simplement due à la variabilité inhérente à l'échantillonnage~?

%%============================================================
\section{Tableau de contingence}
%%============================================================

Lorsqu'on observe deux variables catégoriques sur les mêmes individus, on résume les données
dans un \textbf{tableau de contingence} (parfois appelé tableau croisé).

\textbf{Définition.} Un \textbf{tableau de contingence} à $r$ lignes et $c$ colonnes est un
tableau dont les lignes représentent les $r$~modalités d'une variable~$X$ et les colonnes
représentent les $c$~modalités d'une variable~$Y$. Chaque cellule~$(i,j)$ contient la
fréquence observée $O_{ij}$, c'est-à-dire le nombre d'individus pour lesquels $X = i$ et
$Y = j$.

%%------------------------------------------------------------
\subsection{Fréquences marginales}
%%------------------------------------------------------------

Les \textbf{fréquences marginales} sont les totaux par ligne et par colonne~:
\begin{align*}
  n_{i\cdot} &= \sum_{j=1}^{c} O_{ij}
    && \text{(total de la ligne } i\text{),}\\
  n_{\cdot j} &= \sum_{i=1}^{r} O_{ij}
    && \text{(total de la colonne } j\text{),}\\
  n &= \sum_{i=1}^{r} \sum_{j=1}^{c} O_{ij}
    && \text{(total général).}
\end{align*}

%%------------------------------------------------------------
\subsection{Proportions conditionnelles}
%%------------------------------------------------------------

Les \textbf{proportions conditionnelles} permettent d'examiner informellement la présence
d'une association. La proportion conditionnelle de $Y = j$ sachant $X = i$ est~:
\[
  \hat{p}_{j \mid i} = \frac{O_{ij}}{n_{i\cdot}}.
\]
Si les proportions conditionnelles sont semblables d'une ligne à l'autre (resp. d'une colonne
à l'autre), les données semblent compatibles avec une situation d'indépendance. En revanche,
si elles varient considérablement, cela suggère une association entre $X$ et $Y$.

%%------------------------------------------------------------
\subsection{Exemple illustratif}
%%------------------------------------------------------------

\textbf{Exemple.} On interroge 200 personnes sur leur boisson chaude préférée (Café, Thé,
Chocolat chaud) selon la saison de l'entrevue (Hiver, Été, Automne). Le tableau de contingence
est le suivant~:

\begin{center}
\renewcommand{\arraystretch}{1.6}
\begin{tabular}{|P{3cm}|P{2cm}|P{2cm}|P{2cm}|P{2cm}|}
\hline
\diagbox[innerwidth=3cm]{Boisson}{Saison} & Hiver & Été & Automne & \textbf{Total} \\
\hline
Café          & 40 & 30 & 25 & \textbf{95} \\
\hline
Thé           & 20 & 35 & 15 & \textbf{70} \\
\hline
Chocolat chaud & 25 &  5 &  5 & \textbf{35} \\
\hline
\textbf{Total} & \textbf{85} & \textbf{70} & \textbf{45} & \textbf{200} \\
\hline
\end{tabular}
\end{center}

Par exemple, la proportion conditionnelle de personnes qui choisissent le chocolat chaud
\emph{en hiver} est $25/85 \approx 29{,}4\,\%$, alors qu'elle n'est que $5/70 \approx 7{,}1\,\%$
\emph{en été}. Cette variation suggère qu'il pourrait exister une association entre la boisson
choisie et la saison.

%%============================================================
\section{Indépendance de deux variables catégoriques}
%%============================================================

\textbf{Définition.} Deux variables catégoriques $X$ et $Y$ sont \textbf{indépendantes} si la
distribution conditionnelle de $Y$ est identique pour toutes les modalités de $X$
(et réciproquement). Formellement~:
\[
  P(X = i,\; Y = j) = P(X = i)\cdot P(Y = j)
  \quad \text{pour tout } i \in \{1,\ldots,r\},\; j \in \{1,\ldots,c\}.
\]

Si $X$ et $Y$ sont indépendantes, la fréquence \emph{attendue} dans la cellule $(i,j)$ d'un
tableau de contingence issu d'un échantillon de taille $n$ est~:
\[
  E_{ij}
  = n \cdot P(X = i) \cdot P(Y = j)
  \;\approx\; n \cdot \frac{n_{i\cdot}}{n} \cdot \frac{n_{\cdot j}}{n}
  = \frac{n_{i\cdot} \cdot n_{\cdot j}}{n}.
\]

%%============================================================
\section{Test d'indépendance du khi-carré}
%%============================================================

%%------------------------------------------------------------
\subsection{Hypothèses}
%%------------------------------------------------------------

\[
  H_0 : X \text{ et } Y \text{ sont indépendantes}
  \qquad\text{contre}\qquad
  H_1 : X \text{ et } Y \text{ ne sont pas indépendantes.}
\]

%%------------------------------------------------------------
\subsection{Fréquences attendues sous $H_0$}
%%------------------------------------------------------------

Sous $H_0$, la fréquence attendue dans la cellule $(i,j)$ est estimée par~:
\[
  E_{ij} = \frac{n_{i\cdot} \cdot n_{\cdot j}}{n}.
\]

\remarque{La somme de toutes les fréquences attendues est égale à la somme de toutes les
fréquences observées~: $\sum_{i,j} E_{ij} = n$.}

%%------------------------------------------------------------
\subsection{Statistique de test}
%%------------------------------------------------------------

La statistique du khi-carré mesure l'écart global entre les fréquences observées et les
fréquences attendues sous $H_0$~:
\[
  \chi^2 = \sum_{i=1}^{r} \sum_{j=1}^{c} \frac{(O_{ij} - E_{ij})^2}{E_{ij}}.
\]

\begin{itemize}[itemsep=2pt]
  \item Si $\chi^2 \approx 0$, les fréquences observées sont très proches de celles attendues
        sous l'indépendance~; les données sont compatibles avec $H_0$.
  \item Plus $\chi^2$ est grand, plus les données s'éloignent de l'indépendance, et plus
        l'évidence contre $H_0$ est forte.
\end{itemize}

%%------------------------------------------------------------
\subsection{Distribution sous $H_0$ et degrés de liberté}
%%------------------------------------------------------------

Sous $H_0$, lorsque les conditions de validité sont satisfaites, la statistique $\chi^2$ suit
approximativement une loi du khi-carré à $(r-1)(c-1)$ degrés de liberté~:
\[
  \chi^2 \;\overset{H_0}{\sim}\; \chi^2_{(r-1)(c-1)}.
\]

\textbf{Définition.} Les \textbf{degrés de liberté} du test du khi-carré d'indépendance pour
un tableau $r \times c$ sont $(r-1)(c-1)$.

\textbf{Exemple.} Pour un tableau $3 \times 3$~: $(3-1)(3-1) = 4$ degrés de liberté.
Pour un tableau $2 \times 2$~: $(2-1)(2-1) = 1$ degré de liberté.

%%------------------------------------------------------------
\subsection{Condition de validité}
%%------------------------------------------------------------

L'approximation par la loi du khi-carré est valide lorsque \textbf{toutes} les fréquences
attendues $E_{ij}$ sont supérieures ou égales à~5.

Si cette condition n'est pas satisfaite, on peut~:
\begin{itemize}[itemsep=2pt]
  \item regrouper certaines modalités pour augmenter les fréquences attendues~;
  \item utiliser le test exact de Fisher (applicable aux tableaux $2 \times 2$).
\end{itemize}

%%------------------------------------------------------------
\subsection{Règle de décision}
%%------------------------------------------------------------

On rejette $H_0$ au seuil $\alpha$ si~:
\[
  \chi^2_{obs} > \chi^2_{\alpha;\,(r-1)(c-1)},
\]
où $\chi^2_{\alpha;\,(r-1)(c-1)}$ est le quantile d'ordre $1-\alpha$ de la loi du khi-carré
à $(r-1)(c-1)$ degrés de liberté.

\remarque{Le test du khi-carré d'indépendance est \textbf{toujours unilatéral à droite}. En
effet, lorsque $H_1$ est vraie, les fréquences observées s'éloignent des fréquences attendues,
ce qui fait augmenter la valeur de $\chi^2$. Contrairement aux tests sur la moyenne ou la
proportion, il n'est pas possible de définir une direction précise pour l'hypothèse
alternative~: on dit simplement qu'il \emph{existe} une association.}

%%------------------------------------------------------------
\subsection{Procédure complète -- exemple}
%%------------------------------------------------------------

\textbf{Exemple.} On reprend le tableau de l'exemple de la section~2 (boissons et saisons,
$n = 200$) et on teste formellement l'indépendance au seuil $\alpha = 5\,\%$.

\medskip
\textbf{Étape 1 -- Hypothèses~:}
\begin{itemize}[itemsep=2pt]
  \item $H_0$~: le choix de boisson chaude et la saison sont indépendants.
  \item $H_1$~: il existe une association entre le choix de boisson chaude et la saison.
\end{itemize}

\textbf{Étape 2 -- Fréquences attendues sous $H_0$~:}
\[
  E_{ij} = \frac{n_{i\cdot} \cdot n_{\cdot j}}{n}.
\]

\begin{center}
\renewcommand{\arraystretch}{1.6}
\begin{tabular}{|P{3cm}|P{3.5cm}|P{3.5cm}|P{3.5cm}|}
\hline
\diagbox[innerwidth=3cm]{Boisson}{Saison} & Hiver & Été & Automne \\
\hline
Café
  & $\dfrac{95 \times 85}{200} = 40{,}375$
  & $\dfrac{95 \times 70}{200} = 33{,}25$
  & $\dfrac{95 \times 45}{200} = 21{,}375$ \\[4pt]
\hline
Thé
  & $\dfrac{70 \times 85}{200} = 29{,}75$
  & $\dfrac{70 \times 70}{200} = 24{,}50$
  & $\dfrac{70 \times 45}{200} = 15{,}75$ \\[4pt]
\hline
Chocolat chaud
  & $\dfrac{35 \times 85}{200} = 14{,}875$
  & $\dfrac{35 \times 70}{200} = 12{,}25$
  & $\dfrac{35 \times 45}{200} = 7{,}875$ \\[4pt]
\hline
\end{tabular}
\end{center}

Toutes les fréquences attendues sont supérieures à~5~: la condition de validité est satisfaite.

\textbf{Étape 3 -- Calcul de la statistique de test~:}
\begin{align*}
  \chi^2_{obs}
  &= \frac{(40 - 40{,}375)^2}{40{,}375}
   + \frac{(30 - 33{,}25)^2}{33{,}25}
   + \frac{(25 - 21{,}375)^2}{21{,}375} \\[4pt]
  &\quad + \frac{(20 - 29{,}75)^2}{29{,}75}
   + \frac{(35 - 24{,}50)^2}{24{,}50}
   + \frac{(15 - 15{,}75)^2}{15{,}75} \\[4pt]
  &\quad + \frac{(25 - 14{,}875)^2}{14{,}875}
   + \frac{(5 - 12{,}25)^2}{12{,}25}
   + \frac{(5 - 7{,}875)^2}{7{,}875} \\[4pt]
  &\approx 0{,}003 + 0{,}318 + 0{,}614 + 3{,}193 + 4{,}500 + 0{,}036
           + 6{,}893 + 4{,}286 + 1{,}049 \\[4pt]
  &\approx 20{,}89.
\end{align*}

\textbf{Étape 4 -- Degrés de liberté~:} $(r-1)(c-1) = (3-1)(3-1) = 4$.

\textbf{Étape 5 -- Valeur critique~:} $\chi^2_{0{,}05;\,4} = 9{,}488$.

\textbf{Étape 6 -- Décision~:} Comme $\chi^2_{obs} = 20{,}89 > 9{,}488$, on \textbf{rejette
$H_0$} au seuil de $5\,\%$.

\textbf{Conclusion~:} Il y a une association statistiquement significative entre le choix de
boisson chaude et la saison ($p < 0{,}05$).

%%============================================================
\section{Remarques complémentaires}
%%============================================================

%%------------------------------------------------------------
\subsection{Le test du khi-carré ne mesure pas la force de l'association}
%%------------------------------------------------------------

Le test du khi-carré détermine si une association est \emph{statistiquement significative}
(c'est-à-dire peu probable sous $H_0$), mais il n'indique pas la \emph{force} de cette
association. En particulier, un très grand échantillon peut conduire à rejeter $H_0$ même
pour une association très faible et sans intérêt pratique.

Pour quantifier la force de l'association, on peut calculer le \textbf{V de Cramér}~:
\[
  V = \sqrt{\frac{\chi^2}{n \cdot \min(r-1,\; c-1)}},
\]
dont les valeurs se situent entre~0 (aucune association) et~1 (association maximale).

%%------------------------------------------------------------
\subsection{Lien avec le test de différence de deux proportions}
%%------------------------------------------------------------

Dans le cas particulier d'un tableau $2 \times 2$, le test du khi-carré est équivalent au
test bilatéral de différence de deux proportions. On obtient la relation~:
\[
  \chi^2_{obs} = Z_{obs}^2,
\]
où $Z_{obs}$ est la statistique du test $Z$ pour la différence de deux proportions, et la loi
$\chi^2_1$ est précisément la loi de $Z^2$ lorsque $Z \sim \mathcal{N}(0,1)$.

%%------------------------------------------------------------
\subsection{Résumé de la procédure}
%%------------------------------------------------------------

\begin{center}
\renewcommand{\arraystretch}{1.5}
\begin{tabular}{p{4cm}p{9cm}}
\toprule
\textbf{Étape} & \textbf{Description} \\
\midrule
1. Hypothèses
  & $H_0$~: indépendance \quad vs \quad $H_1$~: association. \\
2. Fréquences attendues
  & $E_{ij} = n_{i\cdot} \cdot n_{\cdot j} / n$ pour chaque cellule. \\
3. Condition de validité
  & Vérifier que $E_{ij} \geq 5$ pour toutes les cellules. \\
4. Statistique
  & $\chi^2 = \displaystyle\sum_{i,j} (O_{ij} - E_{ij})^2 / E_{ij}$. \\
5. Degrés de liberté
  & $\nu = (r-1)(c-1)$. \\
6. Valeur critique
  & $\chi^2_{\alpha;\,\nu}$ (unilatéral à droite). \\
7. Décision
  & Rejeter $H_0$ si $\chi^2_{obs} > \chi^2_{\alpha;\,\nu}$. \\
\bottomrule
\end{tabular}
\end{center}

\end{document}
